\section{Parameters in Gorelov's classification}\label{app:gorelov}
Gorelov uses data $(l,q,p;\boldsymbol\nu,\boldsymbol\lambda,\alpha)$.
His 2021 Lemmas~12 and~14 treat $p_1=p_2=1$ and arbitrary powers,
respectively (statements p.~7; proofs p.~7 and \mbox{pp.}~8--9)
\cite[Lemmas~12 and~14]{Gorelov2021}.
The proof of Lemma~12 in \cite[p.~7]{Gorelov2021} invokes
Lemma~14 of that article's reference~[5] (Gorelov, \emph{Math. Notes}
94 (2013), 82--95), then Lemmas~1--2 of the citing article.
His 2020 Theorem~2 gives mixed sufficiency
\cite[p.~138]{Gorelov2020}, extending the 2019 identity
\cite[(3), p.~260]{Gorelov2019}. The 2019 matrix form
\cite[Theorem~1, p.~260]{Gorelov2019} assumes $2\lambda\notin\Z$
and excludes our case $\nu_0=0$.
Theorem~\ref{thm:linkage} uses
\begin{equation*}
\begin{gathered}
\begin{array}{c|ccc|c|c|c}
 &l&q&p&\boldsymbol\nu&\boldsymbol\lambda&\alpha\\ \hline
 \operatorname{Bess}(c,\nu_0)&0&2&2&\varnothing&(1,1+\nu_0)&c^2/4\\
 \operatorname{Kum}(s,a)&1&2&1&(a)&(1,1)&s
\end{array}
\end{gathered}
\end{equation*}
BBH Theorem~2.2(iv), p.~288, applies to Lemma~\ref{lem:groups}'s scalar
extensions of $\SL_2$ \cite{BBH1988}. It supplies irreducible Siegel-normal
homogeneous systems for Proposition~2.4, \mbox{pp.}~291--292, including logarithmic
second solutions but excluding separately adjoined $L$ and E-parts;
see Supplementary Section~\ref{sup:classical-comparison}. Admissibility conditions: \mbox{pp.}~280--281, 298--299;
rational-value conclusion: Corollary~4.6, \mbox{pp.}~300--301.
Gorelov records the Bessel--Kummer omission and corrected homogeneous condition
\cite[Example~1 and (9), p.~3; Lemma~14, \mbox{pp.}~7--9]{Gorelov2021}.

Scales are nonzero. Denominator parameters avoid $\Z_{\le0}$. Lemma~\ref{lem:groups} gives irreducibility. The formal data are
\begin{equation*}
\begin{gathered}
\begin{array}{c|c|c|c}
 \text{block}&\text{Wronskian, up to a constant}&
   (\text{rate},\text{power})\text{ at }\infty&\text{exponents at }0\\ \hline
 \operatorname{Bess}(c,0)&z^{-1}&(\pm c,-1/2)&0,0\text{ with log}\\
 \operatorname{Bess}(c,\nu)&z^{-2\nu-1}&(\pm c,-\nu-1/2)&0,-2\nu\\
 \operatorname{Kum}(s,a)&e^{sz}/z&(0,-a),(s,a-1)&0,0\text{ with log}
\end{array}
\end{gathered}
\end{equation*}
The powers follow from \cite[10.40.1, 13.7.2--13.7.3]{DLMF} and are
preserved modulo $\Z$ by Lemma~\ref{lem:formal}. At zero, distinct
nonresonant eigenvalues versus nontrivial resonant unipotence give
Theorem~\ref{thm:linkage}'s exclusions.

For the Bessel--Bessel necessity, use the convention
$M_j\simeq M_i\otimes e^{\lambda z}z^\upsilon$.
The rates at infinity give
$\{c',-c'\}=\{c+\lambda,-c+\lambda\}$. Thus $\lambda=0$ and
$c'=\pm c$. The powers then give $\upsilon\equiv\nu_0-\nu'_0\pmod\Z$.
At zero, matching the exponent multisets requires
$\{0,-2\nu'_0\}\equiv\{\upsilon,\upsilon-2\nu_0\}\pmod\Z$.
If $\upsilon\equiv0$, the infinity congruence gives $\nu'_0\equiv\nu_0$. If $\upsilon\equiv2\nu_0$, it gives $\nu'_0\equiv-\nu_0$.
In $[0,1/2)$ the first implies equality and the second forces both zero. Hence $\nu'_0=\nu_0$.
